\chapter{Schädliche Freigebigkeit}
\quellseite{10}

„Selig sind, die Gottes Wort hören und bewahren“, Luc. 11, 28. Ein Missionar in Africa predigte einst über die Worte: „Was hülfe es dem Menschen, so er die ganze Welt gewönne und nähme doch Schaden an seiner Seele?“ Matth. 16, 26., und sagte unter anderm: „Mancher verliert seine Seele dadurch, daß er zu freigebig ist.“ Als er nun sah, daß die Versammlung darüber erstaunt war, so wiederholte er feierlich diese Worte und erklärte seine Meinung dahin: „Viele besuchen die Kirche, hören die Predigt und theilen sie, wenn sie vorbei ist, unter die Leute aus. Dieser Theil ist für diesen Mann, jener Theil für diese Frau; jene Ermahnungen sind für diese Personen und diese Drohungen für diese Sünder; und so geben sie die ganze Predigt weg und behalten nichts für sich.“ Das ist aber eine Erfahrung, die man nicht allein in Africa, sondern auch in America macht. Was hilft es, am Brunnen den Krug füllen bis obenan, wenn man im Heimgehen das Wasser wieder verschüttet? Merke: Das Wort Gottes, das du hörst, gilt vor allem dir, nicht andern.
